%!TeX root=MemoriaTFG.tex

\chapter{Limitaciones}
\label{cap:limitaciones}

El estudio presenta limitaciones de cinco órdenes que conviene
declarar explícitamente. La presente sección las desarrolla por
familia, cuantifica su magnitud cuando es posible y acota su
impacto sobre la interpretación de los resultados del
capítulo~\ref{cap:resultados} y de la discusión del
capítulo~\ref{cap:discusion}.

\section{Limitaciones de disponibilidad}
\label{sec:lim-disponibilidad}

\subsection{Modelos sin pesos públicos}
\label{sec:lim-pesos}

Los pesos de DermFM-Zero~\cite{dermfmzero2026} no son públicos en
la fecha de cierre del presente trabajo. Esta restricción impide
su evaluación empírica directa sobre los doce datasets externos
contrastados y restringe su tratamiento al marco conceptual del
capítulo~\ref{cap:sota}. La comparación entre PanDerm Large y
DermFM-Zero, que la publicación original reporta a favor del
segundo, no puede verificarse de forma independiente con los
mismos protocolos aplicados al resto de la comparativa. La
disponibilidad eventual de los pesos liberaría una comparación
directa con los hallazgos cuantitativos de este trabajo, en
particular sobre las cifras de sondeo lineal de la
tabla~\ref{tab:lp-panderm} y sobre el experimento de equidad de
la tabla~\ref{tab:equidad-fototipo}.

Los modelos cerrados accedidos vía API (GPT-4o, GPT-4o-mini,
Gemini~2.5~Pro, Gemini~2.5~Flash) presentan una limitación
estructural: no permiten controlar la versión exacta del modelo en
el momento de la inferencia ni el posprocesado interno aplicado
por el proveedor. Las cifras de la tabla~\ref{tab:llm-comparativa}
son condicionales a la versión disponible en la fecha del
experimento (marzo de 2026) y no son estrictamente reproducibles
ante actualizaciones silenciosas del proveedor. El coste agregado
de la evaluación (\$4{,}07 aproximadamente entre GPT-4o y
GPT-4o-mini, sin desglose público de tarifa por token de entrada y
salida) introduce una segunda dependencia que limita la
replicación por terceros.

\subsection{Conjuntos de preentrenamiento no liberados}
\label{sec:lim-preentrenamiento}

Los conjuntos de preentrenamiento internos del grupo de la
Universidad Monash empleados en PanDerm~\cite{panderm2025} no son
accesibles para terceros. Esta restricción impide reproducir la
fase de preentrenamiento del modelo y limita la verificación a los
pesos liberados. La consecuencia metodológica es que cualquier
afirmación sobre la composición del corpus de preentrenamiento
(\emph{las cuatro modalidades dermatológicas}, \emph{$\sim 2{,}1$
millones de imágenes}) descansa sobre la declaración de los autores
y no sobre auditoría independiente. La misma restricción aplica al
corpus Derm1M de la familia DermLIP, que se distribuye con
\emph{embeddings} agregados pero sin acceso íntegro al material de
imagen original.

\section{Limitaciones de los datos}
\label{sec:lim-datos}

\subsection{Sesgo geográfico y lingüístico}
\label{sec:lim-geografico}

Los doce datasets externos contrastados provienen mayoritariamente
de instituciones de Estados Unidos (DDI), Australia
(HAM10000 parcial), Europa central (BCN20000 sobre Hospital Clínic
de Barcelona; HAM10000 parcial; ISIC2018; Dermnet) y Brasil
(PAD-UFES-20). HIBA procede del Hospital Italiano de Buenos Aires
y MSKCC del Memorial Sloan Kettering Cancer Center. La población
hispanohablante europea, con la distribución epidemiológica
documentada en DIADERM~\cite{diaderm2018}, no está representada de
forma específica en los conjuntos de evaluación. La consecuencia
es que las cifras agregadas de la tabla~\ref{tab:lp-panderm} no
son directamente extrapolables al espectro diagnóstico atendido en
el Hospital Universitario Son Llàtzer
(HUSLL)~\cite{taberner2010_motivos} sin validación adicional.

La rama lingüística de los modelos vision-lenguaje contrastivos
(DermLIP, BiomedCLIP) está preentrenada sobre corpus en inglés. La
sensibilidad documentada al tokenizador (secci\'on~\ref{sec:disc-zs},
$48{,}8$\,pp AUROC entre la peor y la mejor configuración) se
limita a las variantes evaluadas en inglés y no caracteriza el
rendimiento sobre \emph{prompts} en español, idioma sobre el que
no se han evaluado por la ausencia de validación de los
tokenizadores con vocabulario clínico hispanohablante. El corpus
DermapixelAI~1.0 publicado como contribución del trabajo
(secci\'on~\ref{sec:dermapixel}) aborda esta brecha para la fase
posterior al periodo cubierto por este documento.

\subsection{Heterogeneidad taxonómica entre datasets}
\label{sec:lim-taxonomia}

Las taxonomías diagnósticas nativas de los doce datasets son
incompatibles entre sí. La formulación binaria
melanoma/no-melanoma de Derm7pt clínico, Derm7pt dermo, HIBA y
MSKCC contrasta con la multiclase de HAM10000 ($7$ clases),
PAD-UFES-20 ($6$), DDI ($2$), Dermnet ($23$), BCN20000
($9$), WSI~patches ($16$) y Fitzpatrick17k
($114$ patologías). SkinCon emplea anotación basada en
$48$ conceptos clínicos visuales, taxonomía incompatible con la
diagnóstica de los demás datasets.

La ontología jerárquica L1/L2/L3 introducida en
secci\'on~\ref{sec:ontologia} mitiga esta limitación al armonizar el
dataset agregado sobre vocabulario común
($72\,654$ imágenes etiquetadas, $4$ categorías L1, $43$
subcategorías L2 con cinco entradas sin desagregación a L3, y
$367$ diagnósticos L3 con codificación SNOMED~CT y CIE-10), pero
no la elimina por completo. Las comparaciones inter-dataset son
defendibles a nivel L1 (cuatro familias etiológicas) y L2
($43$ entradas), pero a nivel L3 cada dataset conserva su cola
larga propia y los contrastes directos requieren agregación
adicional. Once de los doce datasets se mapean a la ontología;
SkinCon queda fuera por el desajuste de granularidad ya
mencionado.

La consecuencia operativa para los hallazgos del
capítulo~\ref{cap:resultados} es que el reporte por dataset de la
tabla~\ref{tab:lp-panderm} mantiene fidelidad a la taxonomía
nativa, y las cifras agregadas (filas \emph{Media} de la misma
tabla) deben interpretarse como promedios sobre tareas heterogéneas
no estrictamente comparables entre sí.

\subsection{Solapamiento con el dataset de preentrenamiento Derm1M}
\label{sec:lim-leakage}

La auditoría descrita en secci\'on~\ref{sec:auditoria} identifica tres
datasets evaluados (Dermnet, HIBA, MSKCC) con solapamiento
documentado frente al corpus de preentrenamiento Derm1M
empleado por la familia DermLIP. La cardinalidad exacta de la
intersección por hash MD5 sobre cada uno de los tres datasets es
pendiente de cierre experimental
(tabla~\ref{tab:datasets-resumen}, marcadores con asterisco).

Las métricas obtenidas sobre estos datasets con modelos cuya rama
visual deriva de Derm1M (DermLIP v1, v2 y original) deben
interpretarse con \emph{caveat} y no son directamente comparables
con las obtenidas sobre los datasets sin solapamiento documentado.
La sensibilidad de la métrica AUROC a la presencia de imágenes
previamente vistas en preentrenamiento puede inflar al alza las
cifras de la tabla~\ref{tab:zs-binaria} sobre HIBA y MSKCC,
particularmente en el régimen zero-shot.

Las cifras de PanDerm Base y Large, DINOv2 ViT-L/14,
ConvNeXt-Large, EfficientNetV2-Large y SigLIP-Large SO400M no se
ven afectadas por este solapamiento, dado que su preentrenamiento
no incluye Derm1M. Por consiguiente, los hallazgos cuantitativos
sobre la ventaja del preentrenamiento de dominio
(secci\'on~\ref{sec:disc-especializacion}), la eficiencia en etiquetas
(secci\'on~\ref{sec:disc-eficiencia}) y la comparación con CNNs
(tabla~\ref{tab:gap-cnn}) permanecen defendibles incluso bajo el
peor escenario de la auditoría.

\subsection{Tamaños muestrales por subgrupo en el experimento de equidad}
\label{sec:lim-muestrales}

Los tamaños muestrales del experimento de equidad sobre
Fitzpatrick17k son asimétricos entre fototipos cutáneos: $299$
imágenes para FST~I, $470$ para FST~II, $325$ para FST~III,
$261$ para FST~IV, $169$ para FST~V y $73$ para FST~VI
en el conjunto de test ($N$ total $=1\,597$ con fototipo
asignado). El subgrupo FST~VI dispone de aproximadamente la cuarta
parte de los ejemplos del subgrupo FST~I y supone el $4{,}6\%$
del conjunto de test.

La consecuencia estadística es que los intervalos de confianza al
$95\%$ por \emph{bootstrap} sobre las cifras de
BAcc por subgrupo de la tabla~\ref{tab:equidad-fototipo} son más
amplios para FST~V-VI que para FST~I-IV, lo que limita la
precisión cuantitativa del \emph{gap}~I-VI reportado. La
interpretación cualitativa ---rendimiento inferior sobre fototipos
elevados--- es robusta frente a esta asimetría muestral, pero la
cuantificación exacta del \emph{gap} debe acompañarse de los
intervalos de confianza correspondientes en la versión final del
documento. \textbf{[NO ESPECIFICADO: reportar los intervalos de
confianza al~95\,\% por \emph{bootstrap} sobre la BAcc por
subgrupo y la potencia estadística del contraste entre FST~I y
FST~VI.]}

\section{Limitaciones del diseño experimental}
\label{sec:lim-diseno}

\subsection{Variabilidad inter-semilla}
\label{sec:lim-semilla}

Los protocolos de sondeo lineal, eficiencia en etiquetas y
evaluación de equidad por fototipo operan con una única semilla
(\texttt{random\_state}~$=42$); los protocolos de ajuste fino y
segmentación con LoRA emplean también una única configuración de
semillas (\texttt{torch}, \texttt{numpy}, \texttt{random} fijadas a
$0$). La variabilidad inter-semilla no se reporta de forma
sistemática, lo que limita la cuantificación de la incertidumbre
estadística asociada a la inicialización aleatoria del clasificador
lineal y al orden de muestreo durante el entrenamiento.

La consecuencia es que las cifras de las tablas
\ref{tab:lp-panderm}, \ref{tab:ft-panderm}, \ref{tab:label-eff} y
\ref{tab:equidad-fototipo} son estimaciones puntuales sobre una
trayectoria de optimización concreta y no incluyen su variabilidad
esperable bajo distintas inicializaciones. Las heurísticas
operativas extraídas en secci\'on~\ref{sec:disc-lp-vs-ft}
(régimen $N_{\mathrm{train}} \lesssim 500$ para preferencia de
sondeo lineal) son cualitativas y robustas, pero los umbrales
exactos pueden desplazarse bajo otras semillas.

\subsection{Pruebas estadísticas y comparaciones múltiples}
\label{sec:lim-estadistica}

El protocolo declarado en secci\'on~\ref{sec:estadistica} contempla
intervalos de confianza al $95\%$ mediante \emph{bootstrap} sobre
$1\,000$ remuestreos y el test de DeLong~\cite{delong1988} para
contrastes pareados de AUROC. La función \texttt{resample\_metric}
en el código del proyecto soporta el \emph{bootstrap} estratificado
por clase. Sin embargo, las cifras reportadas en el
capítulo~\ref{cap:resultados} se presentan como valores puntuales,
no como pares (estimador, IC), lo que limita la lectura
estadísticamente informada de los contrastes entre modelos. El
test de DeLong y la corrección por comparaciones múltiples (a
elegir entre Bonferroni y Benjamini-Hochberg) están pendientes de
implementación y ejecución sistemática.

\textbf{[NO ESPECIFICADO: completar la ejecución del bootstrap
estratificado sobre las cifras principales de las tablas
\ref{tab:lp-panderm}, \ref{tab:ft-panderm} y
\ref{tab:equidad-fototipo}; implementar el test de DeLong para los
contrastes pareados de AUROC entre PanDerm Large, DermLIP v2,
DINOv2 y SigLIP-Large; aplicar corrección por comparaciones
múltiples sobre el conjunto de contrastes reportados.]}

La consecuencia operativa es que los hallazgos numéricos del
capítulo~\ref{cap:resultados} son defendibles en términos
descriptivos y de magnitud, pero las afirmaciones sobre
\emph{superioridad estadísticamente significativa} de un modelo
sobre otro requieren la cuantificación pendiente. Los patrones
cualitativos identificados en la discusión
(secci\'on~\ref{sec:disc-especializacion}, secci\'on~\ref{sec:disc-eficiencia},
secci\'on~\ref{sec:disc-lp-vs-ft}) son robustos frente a esta limitación
por la magnitud absoluta de los efectos cuantificados.

\subsection{Sensibilidad al \emph{prompt} en zero-shot y LLMs}
\label{sec:lim-prompts}

La evaluación zero-shot multimodal y la comparación con LLMs
multimodales operan sobre formulaciones concretas de \emph{prompts}
que condicionan los resultados reportados. La diferencia de
$48{,}8$\,pp AUROC entre la peor y la mejor configuración de
DermLIP sobre HAM10000 (tabla~\ref{tab:zs-ham}) cuantifica la
magnitud de esta dependencia para la rama contrastiva. La
generalización a otras formulaciones requiere validación
específica.

Para los LLMs multimodales, el \emph{prompt} clínico de
aproximadamente $300$ palabras descrito en secci\'on~\ref{sec:llm}
introduce una formulación fija que no se ha auditado mediante
barrido sistemático de longitud, orden de las clases candidatas o
inclusión de descriptores anatómicos. El sesgo léxico documentado
en secci\'on~\ref{sec:disc-llm} (GPT-4o emite $284$ predicciones de
dermatofibroma frente a $8$ muestras reales) sugiere que el
modelo responde al espacio léxico del \emph{prompt}, lo que
amplifica el efecto de la formulación. Las cifras de la
tabla~\ref{tab:llm-comparativa} son condicionales a este
\emph{prompt} y no caracterizan el rendimiento de los modelos bajo
formulaciones alternativas. Una caracterización sistemática de la
sensibilidad al \emph{prompt} en LLMs queda fuera del alcance del
trabajo y se documenta como línea futura en el
Anexo~\ref{cap:anexog}.

\section{Limitaciones del corpus DermapixelAI 1.0}
\label{sec:lim-dermapixel}

El corpus DermapixelAI~1.0 entregado como contribución del trabajo
presenta cuatro limitaciones específicas declaradas explícitamente
en el Anexo~\ref{cap:anexoc} y en el documento formal de entrega
(Anexo~\ref{cap:anexoe}).

\paragraph{Desbalance estructural por modalidad.}
La distribución por modalidad es $1\,036$ imágenes de fotografía
clínica frente a $49$ imágenes dermatoscópicas, lo que limita
el uso del corpus para entrenamiento o evaluación de tareas
dermatoscópicas. El predominio clínico refleja la finalidad
docente original del archivo Dermapixel, no una decisión de
diseño del trabajo, pero condiciona la utilidad del corpus para
benchmarks comparativos con HAM10000, BCN20000, ISIC2018 u otros
conjuntos dermatoscópicos.

\paragraph{Cobertura ontológica parcial.}
El corpus contiene anotación L3 sobre $250$ diagnósticos
distintos, sobre las $367$ entradas declaradas en la
ontología, lo que corresponde a una cobertura efectiva del
$68{,}1\%$ del vocabulario L3 con al menos una imagen. El
diagnóstico L3 más frecuente concentra aproximadamente el
$32\%$ del corpus, lo que reproduce la cola larga típica del
material clínico. Esta cobertura parcial limita el uso del corpus
para evaluación granular sobre L3 sin agregación a L2.

\paragraph{Tres cifras de validación no intercambiables.}
La distinción entre los tres campos de validación documentados en
el Anexo~\ref{cap:anexoc} es relevante para la interpretación
correcta de las cifras de calidad. El $97{,}93\%$ del corpus
con \texttt{label\_source = ontology} indica mapeo ontológico
validado del vocabulario, no verificación visual caso-a-caso. El
$82{,}96\%$ con \texttt{rosa\_verified} corresponde a la revisión
visual explícita por la Dra.~R.~Taberner. La cifra de validación
textual \texttt{diagnosis\_source = expert\_v3} ($98{,}38\%$) se
refiere a la coherencia del razonamiento textual con el
diagnóstico asignado, sin verificación visual de la imagen
asociada. Las tres cifras deben declararse con precisión en
cualquier comunicación derivada y no son intercambiables.

\paragraph{Splits y tamaños reducidos.}
La partición \emph{case-aware} del corpus
($891 / 157 / 41$ imágenes en \emph{train/val/test}) es
metodológicamente correcta para evitar fuga de información entre
imágenes del mismo caso clínico, pero produce un conjunto de test
reducido ($41$ imágenes) que limita la potencia estadística de
los contrastes realizados sobre él. La explotación experimental
sistemática sobre DermapixelAI~1.0 se ha pospuesto al periodo
posterior al cierre de este documento por esta razón.

\section{Limitaciones de generalización clínica}
\label{sec:lim-clinica}

\subsection{Material de archivo frente a flujo asistencial real}
\label{sec:lim-archivo}

El estudio mide rendimiento sobre fotografías procedentes de
archivos de investigación curados (HAM10000, BCN20000, ISIC2018) o
de archivos clínicos con verificación posterior (HIBA, MSKCC, DDI,
Fitzpatrick17k). El rendimiento sobre material adquirido en el
flujo real de teledermatología, sobre fotografías de atención
primaria con dispositivos heterogéneos (variabilidad de
iluminación, distancia focal, encuadre, calibración cromática) o
sobre material adquirido en consulta presencial sin estandarización
no se aborda en este trabajo. La extrapolación de las cifras de la
tabla~\ref{tab:lp-panderm} a estos escenarios requiere validación
adicional sobre cohorte representativa del entorno de despliegue
previsto, particularmente en el contexto del proyecto institucional
de teledermatología con IB-Salut descrito en
secci\'on~\ref{sec:trabajo-previo}.

\subsection{Ausencia de validación prospectiva}
\label{sec:lim-prospectivo}

El trabajo no aborda la validación clínica prospectiva sobre
cohorte hospitalaria real con protocolo aprobado por Comité Ético
de Investigación Clínica (CEIC). Esta limitación es estructural y
se asume explícitamente como decisión de alcance en
secci\'on~\ref{sec:alcance}. La validación prospectiva del prototipo
DermApIxel descrito en el Anexo~\ref{cap:anexoh} sobre cohorte del
HUSLL requiere prerrequisitos administrativos y temporales
(aprobación CEIC, definición de \emph{endpoints} clínicos
primarios y secundarios, cálculo muestral, contratación de
revisores independientes) que exceden el alcance del Trabajo de
Fin de Grado y se identifican como línea de continuación inmediata
en la secci\'on~\ref{sec:lineas-abiertas} del capítulo~\ref{cap:conclusiones}.

\subsection{Elementos no abordados por decisión de alcance}
\label{sec:lim-fuera}

Cuatro elementos quedan fuera del alcance declarados en
secci\'on~\ref{sec:alcance} y deben reconocerse como limitaciones del
documento aunque sean por decisión deliberada: (i)~el
preentrenamiento desde cero de un modelo fundacional propio, que
requeriría acceso al material de las cuatro modalidades
dermatológicas en volumen comparable al de PanDerm
($\sim 2{,}1$ millones de imágenes); (ii)~la integración
operativa con el sistema PACS y el visor radiológico del HUSLL,
que excede el alcance técnico del trabajo experimental;
(iii)~el reentrenamiento multilingüe de la rama lingüística sobre
corpus en español de gran escala, condicionado a la disponibilidad
futura de material en cantidad suficiente; (iv)~la certificación
CE como producto sanitario (\emph{Software as a Medical Device}),
que requiere un proceso regulatorio y de auditoría independiente
no contemplado.

\section{Síntesis: dominio de validez de los resultados}
\label{sec:lim-sintesis}

El conjunto de limitaciones declarado en este capítulo acota el
dominio de validez de las cifras reportadas en el
capítulo~\ref{cap:resultados} pero no invalida los patrones
cualitativos identificados en la discusión. Cuatro afirmaciones
cualitativas resisten al peor escenario combinado de las
limitaciones anteriores y constituyen la respuesta empírica
defendible del trabajo a la pregunta de investigación formulada
en secci\'on~\ref{sec:pregunta}.

Primero, la ventaja del preentrenamiento sobre el dominio clínico
respecto a encoders generalistas se cuantifica en
$+4{,}1$\,pp accuracy y $+9{,}0$\,pp BAcc de promedio entre
PanDerm Large y PanDerm Base sobre los diez datasets externos,
patrón estable bajo cualquier interpretación de las
limitaciones de muestreo y semilla. Segundo, la saturación
temprana del sondeo lineal sobre encoders preentrenados a gran
escala produce rendimientos competitivos con $1$--$5\%$ del
conjunto de entrenamiento, hallazgo robusto sobre HAM10000 y
PAD-UFES-20. Tercero, la brecha cuantitativa entre encoders
específicos de dominio y LLMs multimodales generalistas sobre
HAM10000 alcanza $43{,}5$\,pp accuracy entre el mejor
especialista ajustado y el mejor LLM \emph{zero-shot}, magnitud
que sobrevive a la sensibilidad al \emph{prompt} documentada y a
las restricciones de modelos cerrados. Cuarto, el \emph{gap}
sistemático por fototipo cutáneo entre encoders, con sentido
peor-en-fototipos-elevados, es coherente con la literatura previa
y se observa en los cinco modelos contrastados pese a las
limitaciones de tamaño muestral en FST~V-VI.

Estas cuatro afirmaciones constituyen el núcleo defendible del
trabajo y se retoman como conclusiones en el
capítulo~\ref{cap:conclusiones}.
